% 第2章 相关工作
\section{相关工作}

本章综述与本研究相关的技术背景和研究现状，包括无本体数据采集方法\cite{chi2024umi,wang2024dexcap}、双臂机器人遥操作技术\cite{zhao2023aloha,gello2023}、模仿学习策略\cite{argall2009survey,osa2018algorithmic}以及数据质量评估方法\cite{deminf2025,gretton2005hsic}四个方面的内容。

\subsection{无本体数据采集方法}

无本体数据采集是具身智能领域近年来的重要研究方向，其核心思想是在不依赖真实机器人硬件的情况下获取高质量的演示数据。

\textbf{（1）基于手持末端的数据采集}

通用操作接口（Universal Manipulation Interface）是这一类采集方案的典型代表。Chi等人\cite{chi2024umi}设计的系统让操作者手持专门的操作端完成演示，通过配置在腕部的广角相机捕捉视觉信息，利用侧向反光镜辅助估计深度，同时结合惯性测量单元记录手部运动。这种架构使得在真实场景中获取操作演示数据无需机器人本体在场，为模仿学习提供了数据来源。

FastUMI在UMI的基础上进行了硬件和算法优化，提高了数据采集设备的可靠性和易用性。本研究采用的FastUMI系统继承了这一设计思路，使用XR机器人进行双臂数据采集，进一步降低了硬件成本。

\textbf{（2）基于动作捕捉的数据采集}

Wang等人\cite{wang2024dexcap}开发的DexCap方案融合了动作捕捉技术，通过数据手套记录手指的精细运动，结合腕部视觉和惯性传感器追踪手臂轨迹。该方案在UMI基础上增设了胸前广角相机，增强了环境感知的稳定性，适用于需要灵巧操作的复杂任务场景。

\textbf{（3）众包数据采集}

Generalist AI团队开发了指套式采集设备，通过大众力量众包采集机器人演示数据。据报道，该方案峰值时每周可获取超过一万小时的人类演示数据。Sunday Robotics也采用类似方案专注于家庭服务机器人领域。这些工作证明无本体采集配合众包模式可以突破传统数据匮乏的瓶颈。

\subsection{双臂机器人遥操作技术}

双臂协调是机器人操作中的重要课题，遥操作技术为双臂数据采集提供了直接手段。

\textbf{（1）主从机械臂遥操作}

经典的主从遥控制架构要求操作者使用与执行端结构相同的主控设备，通过运动学映射实现动作复现。ALOHA系统\cite{zhao2023aloha}验证了这种架构在家庭操作任务中的可行性，GELLO框架\cite{gello2023}则探索了低成本硬件适配方案。这类方案能够采集高精度的运动轨迹，但设备投入较大，需要同时维护两套机械臂系统。

\textbf{（2）VR遥操作}

Open TeleVision系统\cite{cheng2024opentelevision}通过虚拟现实界面提供沉浸式的操作体验，利用手柄或手套追踪人体动作并实时驱动机器人运动。国内也有团队采用类似技术路线构建了大规模操作数据集。这类方案的优势在于无需操作者配备机械臂本体，但存在操作延迟和缺少触觉反馈的局限。

\textbf{（3）无本体双臂采集}

本研究的无本体双臂采集结合了上述两类方法的优点：既无需昂贵的机械臂本体（降低成本），又能获得人直接操作的精细动作（保证数据质量），同时避免了VR系统的延迟问题。

\subsection{模仿学习策略}

模仿学习（Imitation Learning）是机器人从人类演示中学习策略的主要方法。

\textbf{（1）行为克隆}

行为克隆（Behavior Cloning）是最基础的模仿学习方法\cite{argall2009survey,osa2018algorithmic}，通过监督学习直接从演示数据中学习状态到动作的映射。该方法简单有效，但在分布外数据上容易累积误差\cite{osa2018algorithmic}。

\textbf{（2）扩散策略}

Diffusion Policy\cite{chi2023diffusion}将机器人动作生成建模为去噪扩散过程，能够表达多模态分布并生成平滑轨迹。RDP（Reactive Diffusion Policy）\cite{xue2025rdp}采用"慢-快"双策略架构，以视觉主导产生规划，以触觉反馈快速调整接触细节，在需要精细力控的任务上表现优异。

\textbf{（3）视觉-语言-动作模型}

VLA（Vision-Language-Action）模型将预训练的视觉-语言模型接入机器人控制。Google的RT-2\cite{rt22023}利用海量图像-文本训练数据，实现指令驱动的操作和新技能零样本泛化。本研究采用的π0.5模型延续了这一架构，通过大规模预训练提升策略泛化能力。

\subsection{数据质量评估方法}

数据质量直接影响模仿学习的最终性能，近年来研究者提出了多种评估方法。

\textbf{（1）轨迹分析方法}

轨迹相似性分析使用动态时间规整（DTW）距离度量不同轨迹间的相似程度，结合层次聚类可识别数据中的主要模式。该方法能够发现离群演示，评估数据覆盖度。轨迹大数据综述文献讨论了这些方法在机器人数据挖掘中的应用。

\textbf{（2）互信息估计}

状态-动作互信息可用于衡量数据的可预测性和多样性。Robot Data Curation with Mutual Information Estimators\cite{deminf2025}利用互信息估计进行机器人演示数据的策展。希尔伯特-施密特独立性准则（HSIC）为互信息估计提供了计算框架。

\textbf{（3）分布距离度量}

最大均值差异（MMD）和推土机距离（Wasserstein Distance）可用于度量不同批次数据间的分布差异。MMD Imitation Learning\cite{mmdil2013}利用嵌入分布距离刻画演示与策略差异。本研究的轨迹相似性分析和图像-动作依赖性分析借鉴了这些方法。

综上所述，已有研究为本课题提供了重要基础：无本体采集设备为数据规模化提供了可能，VLA模型和扩散策略为策略学习提供了有效工具，数据质量分析方法为评估采集数据提供了手段。本文在这些基础上，针对双臂协调和系统完整性展开深入研究。
